\documentclass{article}
\usepackage[T1]{fontenc}
\usepackage[utf8]{inputenc}
\usepackage[a4paper, margin=3cm]{geometry}
\usepackage{parskip}
\usepackage{xcolor}
\usepackage{minted}
\usepackage{tabularx}
\usepackage{tikzgitgraph}
\usepackage[hidelinks]{hyperref}

\setminted{fontsize=\small, bgcolor=gray!10, breaklines=true}
\usemintedstyle{pastie}

\newcommand{\cmd}[1]{\texttt{\textbackslash#1}}
\newcommand{\opt}[1]{\texttt{#1}}

\definecolor{cmain}{HTML}{2563EB}
\definecolor{cdev}{HTML}{16A34A}

\title{\texttt{tikzgitgraph}}
\author{}
\date{}

\begin{document}
    \maketitle
    \tableofcontents

    \section{Loading the package}

    \begin{minted}{latex}
\usepackage{tikzgitgraph}
    \end{minted}

    The package automatically loads \texttt{tikz} and the \texttt{calc} library.

    \section{The \texttt{gitgraph} environment}

        All drawing commands must appear inside a \texttt{gitgraph} environment, which wraps a \texttt{tikzpicture}.

        \begin{minted}{latex}
\begin{gitgraph}[options]
...
\end{gitgraph}
        \end{minted}

        Options can also be set globally with \cmd{gitgraphsetup\{...\}}.

        \bigskip
        \begin{tabularx}{\textwidth}{llX}
        \hline
        \textbf{Option} & \textbf{Default} & \textbf{Description} \\
        \hline
        \opt{dir}        & \opt{tb}       & Flow direction: \opt{tb} (top to bottom) or \opt{bt} (bottom to top) \\
        \opt{step}       & \opt{0.75cm}   & Vertical distance between consecutive commits \\
        \opt{xstep}      & \opt{1.6cm}    & Horizontal distance per branch offset unit \\
        \opt{branchdir}  & \opt{right}    & Side on which commit labels are placed (\opt{right} or \opt{left}) \\
        \opt{labelsep}   & \opt{0.2cm}    & Gap between the commit dot and its label \\
        \opt{linewidth}  & \opt{1pt}      & Stroke width of all lines \\
        \opt{dotradius}  & \opt{2pt}      & Radius of the commit dot \\
        \opt{round}      & \opt{4pt}      & Corner rounding of elbows \\
        \hline
        \end{tabularx}

    \section{Declaring a repository: \cmd{repo}}

        \begin{minted}{latex}
\repo[options]{branch-name}{coordinate}
        \end{minted}

        Initialises a branch and places its starting cursor at \texttt{coordinate}.
        Must be called before any \cmd{commit} on that branch.
        Branch names must not contain \texttt{-} or \texttt{:}.

        \bigskip
        \begin{tabularx}{\textwidth}{llX}
        \hline
        \textbf{Option} & \textbf{Default} & \textbf{Description} \\
        \hline
        \opt{color}     & \opt{black}  & Line and dot colour for this branch \\
        \opt{main=true} & \opt{false}  & Draws this branch's lines on the top layer,
                                        above all other branches \\
        \opt{name}      & (empty)      & Text label shown at the starting coordinate \\
        \hline
        \end{tabularx}

    \section{Adding a commit: \cmd{commit}}

        \begin{minted}{latex}
\commit{branch-name}[options]{label}
        \end{minted}

        Advances the branch cursor by one \opt{step}, draws a connecting line
        from the previous commit, and places the dot and label.

        \bigskip
        \begin{tabularx}{\textwidth}{llX}
            \hline
            \textbf{Option} & \textbf{Default} & \textbf{Description} \\
            \hline
            \opt{id} & auto & Explicit node identifier, used when referencing this commit
                            in \cmd{branch} or \cmd{merge} \\
            \hline
        \end{tabularx}

        \subsection{Example: a single branch with three commits}

            \begin{minted}{latex}
\begin{gitgraph}
    \repo[main=true, color=blue, name=main]{main}{0,0}
    \commit{main}{commit1}
    \commit{main}{commit2}
    \commit{main}{commit3}
\end{gitgraph}
            \end{minted}

            \begin{center}
                \begin{gitgraph}
                    \repo[main=true, color=cmain, name=main]{main}{0,0}
                        \commit{main}{commit1}
                        \commit{main}{commit2}
                        \commit{main}{commit3}
                \end{gitgraph}
            \end{center}

    \section{Creating a branch: \cmd{branch}}

        \begin{minted}{latex}
\branch{source-branch:id}{new-branch}[options]
        \end{minted}

        Creates a new branch forking from an existing commit.
        The source is written as \texttt{branch-name:commit-id}.
        A fork elbow is drawn automatically from the source commit to the position of the first future commit on the new branch.

        \bigskip
        \begin{tabularx}{\textwidth}{llX}
            \hline
            \textbf{Option} & \textbf{Default} & \textbf{Description} \\
            \hline
            \opt{color}  & \opt{black} & Colour of the new branch \\
            \opt{offset} & \opt{1}     & Horizontal position relative to the source branch,
                                        in multiples of \opt{xstep}.
                                        Negative values place the branch to the left. \\
            \hline
        \end{tabularx}

        \subsection{Example: a branch}

            \begin{minted}{latex}
\begin{gitgraph}
    \repo[main=true, color=blue, name=main]{main}{0,0}
        \commit{main}{commit1}
        \commit{main}{commit2}
        \branch{main:2}{dev}[color=green!60!black]
            \commit{dev}{dev-commit1}
            \commit{dev}{dev-commit2}
\end{gitgraph}
            \end{minted}

            \begin{center}
                \begin{gitgraph}
                    \repo[main=true, color=cmain, name=main]{main}{0,0}
                        \commit{main}{commit1}
                        \commit{main}{commit2}
                        \branch{main:2}{dev}[color=cdev]
                            \commit{dev}{dev-commit1}
                            \commit{dev}{dev-commit2}
                \end{gitgraph}
            \end{center}

    \section{Merging branches: \cmd{merge}}

        \begin{minted}{latex}
\merge{target-branch}[options]{source-branch:id}{label}
        \end{minted}

        Adds a merge commit on \texttt{target-branch} and draws a merge line from the commit identified by \texttt{source-branch:id}.

        \bigskip
        \begin{tabularx}{\textwidth}{llX}
            \hline
            \textbf{Option} & \textbf{Default} & \textbf{Description} \\
            \hline
            \opt{id}           & auto         & Explicit id for the new merge commit \\
            \opt{parentlonger} & \opt{false}  & The source line drops vertically to align with the target.
                                                Use when the parent has more commits than the child. \\
            \opt{childlonger}  & \opt{false}  & The target extends vertically to align with the source.
                                                Use when the child has more commits than the parent. \\
            \hline
        \end{tabularx}

        \subsection{Merge example: equal number of commits}

            Parent and child branch both have two commits between the branch point and the merge. The default elbow routing works well.

            \begin{minted}{latex}
\begin{gitgraph}
    \repo[main=true, color=blue, name=main]{main}{0,0}
        \commit{main}{commit1}
        \commit{main}{commit2}
        \branch{main:2}{dev}[color=green!60!black]
            \commit{dev}{dev-commit1}
            \commit{dev}{dev-commit2}
        \commit{main}{commit3}
        \commit{main}{commit4}
        \merge{main}{dev:2}{Merge dev}
\end{gitgraph}
            \end{minted}

            \begin{center}
                \begin{gitgraph}
                    \repo[main=true, color=cmain, name=main]{main}{0,0}
                        \commit{main}{commit1}
                        \commit{main}{commit2}
                        \branch{main:2}{dev}[color=cdev]
                            \commit{main}{commit3}
                            \commit{main}{commit4}
                        \commit{dev}{dev-commit1}
                        \commit{dev}{dev-commit2}
                        \merge{main}{dev:2}{Merge dev}
                \end{gitgraph}
        \end{center}

        \subsection{Merge example: parent has more commits (\opt{parentlonger})}

            The parent branch has four commits since the branch point, while the child only has two.
            Without \opt{parentlonger} the merge line would be a steep diagonal; \opt{parentlonger} makes the source line drop vertically first, then elbow horizontally into the target.

            \begin{minted}{latex}
\begin{gitgraph}
    \repo[main=true, color=blue, name=main]{main}{0,0}
        \commit{main}{commit1}
        \commit{main}{commit2}
        \branch{main:2}{dev}[color=green!60!black]
            \commit{dev}{dev-commit1}
            \commit{dev}{dev-commit2}
        \commit{main}{commit3}
        \commit{main}{commit4}
        \commit{main}{commit5}
        \commit{main}{commit6}
        \merge{main}[parentlonger]{dev:2}{Merge dev}
\end{gitgraph}
            \end{minted}

            \begin{center}
                \begin{gitgraph}
                \repo[main=true, color=cmain, name=main]{main}{0,0}
                    \commit{main}{commit1}
                    \commit{main}{commit2}
                    \branch{main:2}{dev}[color=cdev]
                        \commit{dev}{dev-commit1}
                        \commit{dev}{dev-commit2}
                    \commit{main}{commit3}
                    \commit{main}{commit4}
                    \commit{main}{commit5}
                    \commit{main}{commit6}
                    \merge{main}[parentlonger]{dev:2}{Merge dev}
                \end{gitgraph}
            \end{center}

        \subsection{Merge example: child has more commits (\opt{childlonger})}

            The child branch has four commits since the branch point, while the parent only has two.
            \opt{childlonger} extends the target branch vertically to align with the source commit, then elbows the merge line in.

            \begin{minted}{latex}
\begin{gitgraph}
    \repo[main=true, color=blue, name=main]{main}{0,0}
        \commit{main}{commit1}
        \commit{main}{commit2}
        \branch{main:2}{dev}[color=green!60!black]
            \commit{dev}{dev-commit1}
            \commit{dev}{dev-commit2}
            \commit{dev}{dev-commit3}
            \commit{dev}{dev-commit4}
        \commit{main}{commit3}
        \commit{main}{commit4}
        \merge{main}[childlonger]{dev:4}{Merge dev}
\end{gitgraph}
            \end{minted}

            \begin{center}
                \begin{gitgraph}
                    \repo[main=true, color=cmain, name=main]{main}{0,0}
                        \commit{main}{commit1}
                        \commit{main}{commit2}
                        \branch{main:2}{dev}[color=cdev]
                            \commit{dev}{dev-commit1}
                            \commit{dev}{dev-commit2}
                            \commit{dev}{dev-commit3}
                            \commit{dev}{dev-commit4}
                        \commit{main}{commit3}
                        \commit{main}{commit4}
                        \merge{main}[childlonger]{dev:4}{Merge dev}
                \end{gitgraph}
            \end{center}

\end{document}
